% Problem 16: Ramanujan Summation and Roots of Unity Filters
\begin{problem}[{\ramanujanfont Ramanujan Summation \& Filters}]
\begin{enumerate}
    \item[\textbf{(a)}] {\ramanujanfont The Indian mathematician Srinivasa Ramanujan} discovered the series for $1/\pi$:
    \[
    \frac{1}{\pi} = \frac{2\sqrt{2}}{9801} \sum_{n=0}^{\infty} u_n(1103 + 26390n), \text{ where } u_n = \frac{(4n)!}{(n!)^4(396)^{4n}}. \text{ Do NOT prove this.}
    \]

    \begin{enumerate}
        \item[(i)] Show that for $n \ge 0$, the ratio of consecutive terms in the sequence $u_n$ is given by:
        \[
        \frac{u_{n+1}}{u_n} = \frac{(4n + 1)(4n + 3)(2n + 1)}{K(n + 1)^3}, \text{ where } K = \frac{396^4}{8}.
        \]

        \item[(ii)] Hence, prove that for all $n \ge 1$:
        \[
        \frac{u_{n+1}}{u_n} < \frac{1}{9 \times 10^9}
        \]

        \item[(iii)] Let $S_k$ be the sum of the first $k$ terms of the series. Using the result from part (ii), explain why the first term alone (i.e., $n = 0$) provides an approximation for $\pi$ that is accurate to at least 6 decimal places.
    \end{enumerate}

    \smallskip
    \item[\textbf{(b)}] {\ramanujanfont Ramanujan discovered distinct patterns in the partition function, $p(n)$, which counts the number of ways an integer $n$ can be written as a sum. For example, he proved that if $n$ is of the form $5k + 4$ (like 4, 9, 14...), then $p(n)$ is always divisible by 5. To isolate these specific terms, mathematicians often use filters.}

    \noindent Let $\omega = e^{i\frac{2\pi}{5}}$ be a complex fifth root of unity.

    \begin{enumerate}
        \item[(i)] Prove the identity:
        \[
        \sum_{r=0}^{4} \omega^{kr} = 
        \begin{cases} 
        5 & \text{if } k \text{ is a multiple of } 5 \\ 
        0 & \text{if } k \text{ is not a multiple of } 5 
        \end{cases}
        \]

        \item[(ii)] Consider the polynomial $P(z) = \sum_{n=0}^{N} a_n z^n$. Using part (i), show that:
        \[
        \sum_{r=0}^{4} \omega^{-4r} P(\omega^r z) = 5 \sum_{k} a_{5k+4} z^{5k+4}
        \]
        \textit{Here, the symbol $\sum_{k}$ means we are summing over all integers $k$ so that $5k+4$ is an integer in the range from $0$ to $N$.}

        \item[(iii)] If $a_n$ represents the partition numbers $p(n)$, and we assume Ramanujan's congruence $p(5k+4) \equiv 0 \pmod 5$ holds, determine the value of:
        \[
        \sum_{r=0}^{4} \omega^{-4r} P(\omega^r)
        \]
        \textit{Here, the symbol $\equiv$ means ``is congruent to'' (acting like an equals sign), and $\pmod 5$ means we are looking at the remainder when divided by 5; together, $\equiv 0 \pmod 5$ simply means ``is divisible by 5.''}
    \end{enumerate}
\end{enumerate}

\end{problem}

\begin{hint}
\begin{itemize}
    \item \textbf{(a)(i)} When substituting $n+1$, remember that $(4(n+1))! = (4n+4)!$, which expands to $(4n+4)(4n+3)(4n+2)(4n+1)(4n)!$. Write out the ratio carefully and cancel common factorial terms.
    \item \textbf{(a)(ii)} The expression for the ratio is strictly decreasing for $n \ge 1$. Evaluate the ratio at its maximum possible value (when $n=1$) and compare it to the given bound.
    \item \textbf{(a)(iii)} If each subsequent term is less than $10^{-9}$ times the previous term, how fast is the error shrinking? Think about geometric progressions and decimal place accuracy.
    \item \textbf{(b)(i)} Recognize this as a finite geometric series with a common ratio of $\omega^k$. Apply the geometric sum formula for the case where $\omega^k \neq 1$.
    \item \textbf{(b)(ii)} Expand $P(\omega^r z)$ into its sum form. Swap the order of summation so you evaluate the sum over $r$ first. Notice that $\omega^{-4r} \omega^{rn} = \omega^{r(n-4)}$. Apply the result from (b)(i).
    \item \textbf{(b)(iii)} Substitute $z=1$ into your result from (b)(ii). Relate $a_n$ to the partition function $p(n)$, and use the fact that $p(5k+4)$ is a multiple of 5.
\end{itemize}
\end{hint}

\begin{solution}
\small
\textbf{(a)(i)} Starting from the definition of $u_n$,
\[
\frac{u_{n+1}}{u_n}
= \frac{(4n+4)!}{((n+1)!)^4 396^{4n+4}} \cdot \frac{(n!)^4 396^{4n}}{(4n)!}
= \frac{(4n+4)(4n+3)(4n+2)(4n+1)}{(n+1)^4 396^4}.
\]
Now $4n+4=4(n+1)$ and $4n+2=2(2n+1)$, so
\[
\frac{u_{n+1}}{u_n}
= \frac{8(n+1)(4n+3)(2n+1)(4n+1)}{(n+1)^4 396^4}
= \frac{(4n+1)(4n+3)(2n+1)}{K(n+1)^3},
\qquad K=\frac{396^4}{8}.
\]

\textbf{(a)(ii)} The ratio decreases for $n\ge1$, so its largest value occurs at $n=1$. Then
\[
\frac{u_2}{u_1}
= \frac{5\cdot7\cdot3}{K\cdot2^3}
= \frac{105}{396^4}.
\]
Since $396^4\approx 2.46\times 10^{10}$, this is about $4.3\times10^{-9}$, so certainly
\[
\frac{u_{n+1}}{u_n}<\frac{1}{9\times10^9}
\qquad (n\ge1).
\]

\textbf{(a)(iii)} By part (ii), each term is less than $10^{-9}$ times the previous one, so the tail is bounded by an infinite geometric series with $a = u_1$ and $r = 10^{-9}$. This series converges to a value far smaller than the $6^{\text{th}}$ decimal place, making the tail completely negligible. In particular, the $n=1$ term affects only about the $9$th or $10$th decimal place. Hence $S_0$ alone already gives at least $6$ correct decimal places for $1/\pi$.

\textbf{(b)(i)} Let $S=\sum_{r=0}^4 (\omega^k)^r$. If $5\mid k$, then $\omega^k=e^{i2\pi m}=1$, so $S=5$. Otherwise $\omega^k\ne1$, and
\[
S=\frac{1-(\omega^k)^5}{1-\omega^k}
=\frac{1-(\omega^5)^k}{1-\omega^k}
=\frac{1-1}{1-\omega^k}=0.
\]

\textbf{(b)(ii)} Since $P(z)=\sum_{n=0}^N a_n z^n$,
\[
\sum_{r=0}^4 \omega^{-4r}P(\omega^r z)
= \sum_{n=0}^N a_n z^n \sum_{r=0}^4 \omega^{r(n-4)}.
\]
By part (i), the inner sum is $5$ when $n=5k+4$ and $0$ otherwise, so
\[
5\sum_k a_{5k+4}z^{5k+4}.
\]

\textbf{(b)(iii)} Setting $z=1$ in part (ii) gives
\[
\sum_{r=0}^4 \omega^{-4r}P(\omega^r)=5\sum_k a_{5k+4}.
\]
If $a_n=p(n)$, this becomes $5\sum_k p(5k+4)$. Since $p(5k+4)=5m$,
\[
5\sum_k p(5k+4)=25\sum_k m,
\]
so the value is a multiple of $25$.
\normalsize
\end{solution}

\begin{takeaways}
\begin{enumerate}
    \item \textbf{Hyper-Fast Convergence:} Ramanujan's series for $1/\pi$ demonstrates that analyzing the ratio of consecutive terms is key to proving how fast an infinite series converges to its target value.
    \item \textbf{Roots of Unity Filter:} This is an elegant and immensely powerful algebraic tool. By summing over roots of unity, you cause unwanted terms to perfectly cancel each other out (summing to 0), cleanly isolating periodic sequences within a polynomial or power series.
    % \item \textbf{Modular Arithmetic in Combinatorics:} Ramanujan's congruences reveal that the partition function $p(n)$, despite simply counting additive combinations, contains deep, predictable algebraic structures when viewed through the lens of modular arithmetic.
    \item \textbf{Connecting Universes:} Factorials, complex numbers, modular arithmetic, and $\pi$ shouldn't logically sit in the same problem. This is a masterclass in how seemingly isolated topics weave together to solve otherwise impossible problems.
    \item \textbf{Roots of unity and filters (standard identity):} For a primitive $n$th root of unity $\omega = e^{2\pi i/n}$, the sum
    \[
      \sum_{r=0}^{n-1} \omega^{kr} =
      \begin{cases}
        n & \text{if } n \mid k,\\
        0 & \text{otherwise,}
      \end{cases}
    \]
    acts like an on/off switch: it is $n$ exactly when the exponent index $k$ is a multiple of $n$, and $0$ otherwise.
    \item \textbf{Mini-example (polynomial filter):} If $P(z) = \sum_{m=0}^{N} a_m z^m$ and $\omega$ is a primitive $n$th root of unity, then
    \[
      \sum_{r=0}^{n-1} P(\omega^r z)
      = \sum_{m=0}^{N} a_m z^m \sum_{r=0}^{n-1} \omega^{rm}
      = n \sum_{k} a_{kn} z^{kn},
    \]
    because the inner sum is $n$ only when $m$ is a multiple of $n$ and $0$ otherwise. In other words, this sum ``filters out'' and keeps only the coefficients whose indices are multiples of $n$—exactly the mechanism used in part (b) to isolate the $5k+4$ terms.
\end{enumerate}
\end{takeaways}
